\clearpage
\chapter*{2026 Septiembre}
\addcontentsline{toc}{chapter}{2026 Septiembre}

\begin{enunciado}
    \hypertarget{2026 S 1}{}
    \textbf{1. (3 puntos)} Sea \((X, T)\), un espacio topológico y \(M\) un subconjunto de \(X\). Si \(M\) es conexo, estudiar si \(\overline{M}\), adherencia de \(M\), es conexo.
\end{enunciado}

\begin{enunciado}
    \hypertarget{2026 S 2}{}
    \textbf{2. (3 puntos)} Se define la topología de los complementos finitos \(T_{CF}\) sobre \(\R\) como sigue: 
    \[
    T_{CF}=\{\emptyset\}\cup\{A\subset R \mid \R-A \text{ es un conjunto finito}\}
    \]
    Sea \(f\) la aplicación de \(\R\) en \((\R,T_{CF})\), definida para todo \(x \in \R\) por \(f(x) = [x]\), siendo \([x]\) la parte entera de \(x\), es decir, \([x] = \max\{n \in \Z \mid n \leq x\}\), y sea \(T\) la topología de \(\R\) inicial para la aplicación \(f\). Estudie si el subconjunto \(D = \{n + \frac{1}{2} \mid n \in \Z\}\) es denso en \((\R,T)\).
\end{enunciado}

\begin{enunciado}
    \hypertarget{2026 S 3}{}
    \textbf{3. (4 puntos)} En el conjunto \(\R\) de los números reales, definimos \(T\) como la topología definida por la subbase \(S = \{[a,b] \mid a, b \in \Q, a < b\}\). Entonces:
    \begin{enumerate}[label=\alph*)]
        \item Encontrar una base de \((\R,T)\).
        \item Estudiar si \((\R,T)\) es \(T_2\) o Hausdorff.
        \item Estudiar si \((\R,T)\) es separable.
        \item Estudiar si \((\R,T)\) es compacto.
    \end{enumerate}
\end{enunciado}